\subsection{Materiales}

Para la obtención de la información se emplearon como técnicas principales la entrevista, la encuesta y la observación directa cronometrada, aplicadas al personal y a los pacientes del Centro Odontológico Bellidental en sus instalaciones ubicadas en el cantón Baños, provincia de Tungurahua.

Una entrevista semiestructurada fue aplicada a las tres personas que intervienen directamente en la atención al paciente: la secretaria, el doctor especialista y el doctor odontólogo general. El guion fue común para los tres, con un bloque adicional exclusivo para los profesionales médicos sobre las exigencias del formulario de historia clínica y la normativa de confidencialidad, aspecto ajeno al personal administrativo. 

Una encuesta cerrada se dirigió exclusivamente a la secretaria, orientada a cuantificar la carga operativa de coordinar citas por canales no integrados: mensajes y llamadas atendidos por día, tiempo dedicado a confirmar las citas del día siguiente, tiempo de respuesta en horas de alta afluencia y frecuencia de citas perdidas o conflictos de horario. 

Adicionalmente, se aplicó una encuesta a los pacientes, orientada a identificar el canal habitual para agendar citas, las dificultades experimentadas con la comunicación actual (olvidos, confusión de horarios, cancelaciones de última hora) y su disposición a interactuar con un asistente automatizado por WhatsApp para gestiones sencillas. Este instrumento se aplicó en dos modalidades para ampliar su alcance: digital, mediante un formulario de Google Forms accesible por código QR disponible en la clínica, y física, mediante una versión impresa. 

Finalmente, se empleó una ficha de observación cronometrada para registrar la duración de cada tramo administrativo de la atención: llenado o actualización de la historia clínica, registro posterior a la consulta, registro del pago y agendamiento de la siguiente cita, junto con referencias de tiempo como la hora de llegada del paciente y el inicio y fin de la atención clínica. La toma de datos inicial se realizó bajo condiciones operativas habituales, y se prevé repetirla tras implementar el sistema, a fin de contar con una medición directa y comparable del tiempo administrativo antes y después de la intervención, en lugar de depender de percepciones subjetivas del personal.